\documentclass[11pt,a4paper]{article}

\usepackage[margin=2.5cm]{geometry}
\usepackage{amsmath}
\usepackage{amssymb}
\usepackage{bm}
\usepackage{braket}
\usepackage{booktabs}
\usepackage{hyperref}
\usepackage{xcolor}

\title{Why a circular-dichroism analysis of a detuned dimer reports an
inflated coupling}
\author{Derivation note}
\date{\today}

\newcommand{\muv}{\bm{\mu}}
\newcommand{\Rv}{\mathbf{R}}

\emergencystretch=1em

\begin{document}
\maketitle

\begin{abstract}
We derive, from the two-site exciton Hamiltonian and the coupled-oscillator
expression for rotational strength, the quantity that a circular-dichroism (CD)
measurement of an inhomogeneous ensemble of dimers actually reports. Each
configuration contributes to the CD couplet with weight proportional to
$|J|/\Omega$, where $\Omega=\sqrt{\Delta^{2}+4J^{2}}$ is its exciton splitting.
In the \emph{resolved} regime (splitting larger than the linewidth) the first
moment of the couplet is a \emph{harmonic}-type mean of $\Omega$, and an
analyst who fits the couplet under the assumption of degeneracy recovers an
apparent coupling
$J_{\mathrm{app}}=\tfrac{1}{2}\langle\Omega\rangle_{\mathrm{CD}}\ge|J|$. In the
\emph{overlapped} regime (splitting smaller than the linewidth, the operative
case for room-temperature protein chromophores) the couplet collapses to the
derivative of the lineshape: the peak-to-peak separation then reports the
\emph{linewidth}, not the splitting, and only the couplet amplitude and the
signed first moment carry coupling information. The signed first moment is
shown to measure $|JT|$ exactly, free of detuning, at any degree of overlap.
We then read the two published reductions of the fluorescent-protein
tandem-dimer couplet against this framework: the reported splitting is smaller
than the width of either band it is meant to separate, the degeneracy
assumption is hard-coded in the constrained fit rather than tested, and the
couplet centre is taken from the one-chromophore control whose absorption
maximum the two-chromophore construct does not share.
Every assumption is stated explicitly, and those that are not automatically
satisfied are flagged. A closing section describes how to compute the CD
spectrum directly from QM/MM simulation, which removes the inversion problem
altogether.
\end{abstract}

\section{The two-site problem}

Take two chromophores $A$ and $B$ with site energies $E_{A}$, $E_{B}$ and
electronic coupling $J$, in the single-excitation subspace of the standard
molecular-exciton model \cite{davydov1962,kasha1965}:
\begin{equation}
H=\begin{pmatrix} E_{A} & J \\ J & E_{B}\end{pmatrix} .
\end{equation}
Write the mean and difference
\begin{equation}
\bar{E}=\tfrac{1}{2}(E_{A}+E_{B}), \qquad \Delta = E_{A}-E_{B} .
\end{equation}
Diagonalising,
\begin{equation}
E_{\pm}=\bar{E}\pm\tfrac{1}{2}\Omega,
\qquad
\boxed{\;\Omega=\sqrt{\Delta^{2}+4J^{2}}\;}
\label{eq:omega}
\end{equation}
so the observable splitting is $\Omega$, and reduces to $2|J|$ only when
$\Delta=0$. Introduce the mixing angle $\theta$ through
\begin{equation}
\sin 2\theta = \frac{2J}{\Omega},
\qquad
\cos 2\theta = \frac{\Delta}{\Omega},
\label{eq:angle}
\end{equation}
which fixes $2\theta$ uniquely (the frequently quoted
$\tan 2\theta = 2J/\Delta$ alone does not: it leaves a branch ambiguity that
scrambles the state labels when $\Delta$ changes sign). The eigenvectors are
\begin{align}
\ket{+} &= \cos\theta\,\ket{A}+\sin\theta\,\ket{B}, \\
\ket{-} &= -\sin\theta\,\ket{A}+\cos\theta\,\ket{B},
\end{align}
and the product of site amplitudes, which is the quantity that will carry the
chirality, is
\begin{equation}
c^{\pm}_{A}c^{\pm}_{B}=\pm\tfrac{1}{2}\sin 2\theta
=\pm\frac{J}{\Omega}.
\label{eq:mixing}
\end{equation}
Equation~(\ref{eq:mixing}) is the crux: it is bounded in magnitude by
$\tfrac{1}{2}$ at degeneracy and decays as $|J|/|\Delta|$ once
$|\Delta|\gg 2|J|$. Note that it carries the sign of $J$ automatically: for
$J<0$ the symmetric combination becomes the \emph{lower} exciton state and the
couplet inverts, as it should.

\section{Rotational strength of the couplet}

In the coupled-oscillator (exciton-chirality) approximation
\cite{tinoco1962,schellman1975,harada1983,berova2007}, the rotational
strength of the exciton band $k=\pm$ of a dimer is
\begin{equation}
R_{\pm} = -\,\pi\bar\nu\;
c^{\pm}_{A}c^{\pm}_{B}\;
\Rv_{AB}\cdot(\muv_{A}\times\muv_{B}),
\label{eq:rgeneral}
\end{equation}
where $\Rv_{AB}=\Rv_{B}-\Rv_{A}$ is the inter-chromophore displacement,
$\muv_{A},\muv_{B}$ are the monomer electric transition dipoles, and $\bar\nu$
is the mean transition wavenumber.\footnote{Two conventions are folded into the
overall sign of Eq.~(\ref{eq:rgeneral}): the direction of $\Rv_{AB}$ and the
sign convention for the magnetic transition-dipole operator. Only the
\emph{relative} sign of the two bands and their common magnitude matter for
everything below. Using the band-specific $\bar\nu_{\pm}$ instead of a common
$\bar\nu$ changes $R_{\pm}$ by $O(\Omega/\bar\nu)$ (well under $1\%$ here) and
slightly breaks exact conservativity; the common-$\bar\nu$ form is the standard
choice.} The sign of the mixing product must \emph{not} be double-counted: the
familiar degenerate-dimer formula
$R_{\pm}=\mp(\pi\bar\nu/2)\,\Rv_{AB}\cdot(\muv_{A}\times\muv_{B})$ of the
exciton-chirality literature \cite{harada1983,berova2007} is recovered from
Eq.~(\ref{eq:rgeneral}) by inserting $c^{\pm}_{A}c^{\pm}_{B}=\pm\tfrac12$; the
$\mp$ there is nothing but the sign of the amplitude product.\footnote{An
earlier draft of this note wrote
$R_{\pm}=\mp(\pi\bar\nu/2)(2c^{\pm}_{A}c^{\pm}_{B})\,T$, i.e.\ kept the $\mp$
\emph{and} the signed product. Substituting Eq.~(\ref{eq:mixing}) then gives
both bands the same sign, $R_{+}=R_{-}$, contradicting the conservativity used
three lines later. The form (\ref{eq:rgeneral}) is the consistent one.}
Substituting Eq.~(\ref{eq:mixing}) and abbreviating the chirality triple
product as
\begin{equation}
T \equiv \Rv_{AB}\cdot(\muv_{A}\times\muv_{B}),
\end{equation}
gives
\begin{equation}
\boxed{\;
R_{\pm} = \mp\,\pi\bar\nu\,\frac{J}{\Omega}\,T \;}
\label{eq:rotational}
\end{equation}
The couplet is conservative, $R_{+}=-R_{-}$, and its amplitude is controlled by
$|J|/\Omega$. Detuning does not merely widen the couplet; it \emph{weakens} it,
and does so in inverse proportion to the very splitting it creates. The product
$|R_{\pm}|\,\Omega=\pi\bar\nu\,|J|\,|T|$ is independent of $\Delta$, a sum rule
we use repeatedly below.

\section{Ensemble average}

A room-temperature ensemble (or an MD trajectory) samples configurations
$i=1\ldots N$ with different $\Delta_{i}$, and in general different $T_{i}$ and
$J_{i}$. Assume each configuration contributes an identical normalised
lineshape $g$ centred on its own exciton energies, so that the spectrum is the
incoherent sum\footnote{Strictly $\Delta\varepsilon(\nu)\propto
\nu\sum_i[\ldots]$; the frequency prefactor shifts moments by
$O(\sigma_g^{2}/\bar\nu)$ and is suppressed throughout.}
\begin{equation}
\Delta A(\nu)\;\propto\;\frac{1}{N}\sum_{i}\Big[
R^{i}_{+}\,g\!\left(\nu-\nu^{i}_{+}\right)
+R^{i}_{-}\,g\!\left(\nu-\nu^{i}_{-}\right)\Big],
\label{eq:ensemble}
\end{equation}
and assume further that $\mathrm{sign}(J_{i}T_{i})$ is the same for every
configuration, so that all couplets add with the same sense rather than
partially cancelling (see assumption~2 below).

\subsection{Resolved regime}

Suppose first that the couplet is \emph{resolved}: the typical splitting
exceeds the total width of a single-configuration band, so the positive-going
lobe of the measured spectrum is built from the $+$ components only (and
likewise for the negative lobe). Because $g$ is common to all configurations
and normalised, the first moment of each lobe is then the $|R|$-weighted mean
of its component positions, and the separation of the two lobe moments is
\begin{equation}
\langle\Omega\rangle_{\mathrm{CD}}
=\frac{\sum_{i}|R^{i}|\;\Omega_{i}}{\sum_{i}|R^{i}|}
=\frac{\sum_{i}\left(|J_{i}T_{i}|/\Omega_{i}\right)\Omega_{i}}
       {\sum_{i}\left(|J_{i}T_{i}|/\Omega_{i}\right)}
=\frac{\sum_{i} |J_{i}T_{i}|}{\sum_{i} |J_{i}T_{i}|/\Omega_{i}} .
\label{eq:weighted}
\end{equation}
(The separation of lobe moments is insensitive to disorder in $\bar\nu_i$,
since $\nu^{i}_{+}-\nu^{i}_{-}=\Omega_i$ configuration by configuration.) If
$J_{i}$ and $T_{i}$ are constant across the ensemble this collapses to
\begin{equation}
\langle\Omega\rangle_{\mathrm{CD}}
=\frac{N}{\sum_{i}\Omega_{i}^{-1}}
\equiv \Omega_{\mathrm{harm}},
\label{eq:harmonic}
\end{equation}
the \emph{harmonic} mean of the exciton splitting. Since the harmonic mean of a
positive spread quantity never exceeds its arithmetic mean, with equality only
for a degenerate distribution,
\begin{equation}
\langle\Omega\rangle_{\mathrm{CD}}\;\le\;\langle\Omega\rangle ,
\end{equation}
CD systematically reports \emph{less} than the typical splitting of a
disordered ensemble: the $1/\Omega_i$ weighting means the couplet is dominated
by the near-degenerate sub-ensemble. The bias grows with the width of the
$\Omega$ distribution. The analogous weighting of polarized spectra of
disordered aggregates is analysed in Ref.~\cite{somsen1996}.

\subsection{Overlapped regime: the peak separation measures the linewidth}
\label{sec:overlap}

The identification of lobe moments with component positions fails when the two
bands overlap, and this is not a small correction: it changes what the
experiment measures. For a single configuration with $\Omega\ll\sigma_g$
(lineshape width), expand Eq.~(\ref{eq:ensemble}):
\begin{equation}
\Delta A(\nu)\;\approx\;\mp\,|R|\,\Omega\,g'(\nu-\bar\nu),
\label{eq:derivative}
\end{equation}
the \emph{derivative} of the lineshape. For a Gaussian $g$ of standard
deviation $\sigma_g$ the extrema of $g'$ sit at $\pm\sigma_g$; for a Lorentzian
of half-width at half-maximum $\gamma$ they sit at $\pm\gamma/\sqrt{3}$. The
apparent peak-to-peak separation therefore saturates at
\begin{equation}
\Delta\nu_{\mathrm{pp}}\;\xrightarrow[\;\Omega\ll w\;]{}\;
\begin{cases}
2\sigma_g & \text{(Gaussian)}\\[2pt]
2\gamma/\sqrt{3} & \text{(Lorentzian)}
\end{cases}
\qquad\text{independent of }\Omega ,
\label{eq:pp}
\end{equation}
where $w$ denotes the relevant width. The Lorentzian case is not academic: both
published analyses of these constructs fit Lorentzians
(Sec.~\ref{sec:published}). Meanwhile the couplet amplitude, by the sum rule
$|R|\,\Omega=\pi\bar\nu|J||T|$, becomes
\begin{equation}
|\Delta A|_{\max}\;\propto\;|R|\,\Omega\,|g'(\sigma_g)|
\;\propto\;|J|\,|T|,
\end{equation}
independent of $\Delta$. In the overlapped regime the roles therefore swap:
the \emph{amplitude} carries the coupling and the \emph{peak positions} carry
the linewidth \cite{rodger1997,pescitelli2016}. An analyst who reads the
peak-to-peak separation as $2J$ in this regime recovers
$J_{\mathrm{app}}\approx\sigma_g$ (or $\gamma/\sqrt3$) --- the linewidth in
disguise, with no coupling content at all.

The fluorescent-protein constructs that motivate this note sit in exactly this
regime, and their own published parameters say so. Nguyen \emph{et al.}\ fit
the $\mathrm{dVenus\mbox{-}TD}$ difference-ellipticity couplet with Lorentzians
of half-widths $\gamma_{1}=293.5$ and $\gamma_{2}=439.3~\mathrm{cm^{-1}}$ and a
component separation of $2\delta=261.6~\mathrm{cm^{-1}}$ (Table~S3 of
Ref.~\cite{nguyen2025}): the splitting they report is \emph{smaller than either
band it is meant to separate}. Equation~(\ref{eq:pp}) then puts the
derivative-limit floor on any extremum-based reading at
$2\gamma/\sqrt{3} = 339$--$507~\mathrm{cm^{-1}}$. For comparison, the
peak-to-peak Davydov splitting reported by Kim \emph{et al.}\ for
concentration-driven $\mathrm{Venus_{A206}}$ dimers, $14.6\pm0.3$~nm about
$516$~nm \cite{kim2019}, is $548~\mathrm{cm^{-1}}$, i.e.\ $U=274~\mathrm{cm^{-1}}$
--- inside the band that a couplet of \emph{zero} splitting and these widths
would produce.\footnote{Kim's couplet is measured on free
$\mathrm{Venus_{A206}}$ against monomeric $\mathrm{Venus_{A206K}}$ at three
concentrations, not on the tethered tandem construct; the two experiments are
on different species and only Nguyen's is a tandem dimer.}

\subsection{What the published fits actually constrain}
\label{sec:published}

It is worth being precise about the procedure, because the two published
reductions differ from each other and neither is a naive reading of the couplet
extrema. Nguyen \emph{et al.}\ report two \cite{nguyen2025}:

\begin{description}
\item[Unconstrained (main text).] A sum of three Lorentzians in
\emph{wavelength}, $\sum_i A_i/[(\lambda-\lambda_i)^{2}+B_i]$, with all centres
free. The couplet is components $2$ and $3$, and the coupling is defined as
half their wavenumber separation,
$U=\tfrac{1}{2}(\lambda_{2}^{-1}-\lambda_{3}^{-1})\times10^{7}$. For
$\mathrm{dVenus\mbox{-}TD}$, $\lambda=511.8$ and $521.7$~nm give
$U=185.4~\mathrm{cm^{-1}}$.
\item[Constrained (Table~S3).] A sum of three Lorentzians in \emph{wavenumber}
in which the two couplet centres are pinned symmetrically at $G_{1}\mp\delta$,
with $G_{1}$ \emph{fixed} to the first Gaussian of the corresponding
$\mathrm{TDX}$ \emph{absorption} spectrum and $\delta$ the single free
splitting parameter. Widths and amplitudes are independent, and the third band
is free. For $\mathrm{dVenus\mbox{-}TD}$, $G_{1}=19{,}322~\mathrm{cm^{-1}}$ and
$\delta=130.79~\mathrm{cm^{-1}}$; this is the origin of the value usually
quoted.
\end{description}

Four features of the constrained fit bear directly on the derivation above.

\emph{(i) The degeneracy assumption is hard-coded, not implicit.} Placing the
two components at $G_{1}\mp\delta$ and identifying $\delta$ with the coupling
is Eq.~(\ref{eq:omega}) evaluated at $\Delta=0$. The supplement states the
motivation as ``the predicted symmetry expected for Davydov splitting around
this absorption peak when plotted on an energy scale.'' Nothing in the data
tests it; $\Delta$ is never a parameter.

\emph{(ii) The centre is taken from the wrong construct.} $G_{1}$ is the
$\mathrm{TDX}$ absorption peak, but the same analysis (Table~S2 of
Ref.~\cite{nguyen2025}) reports that $G_{1}$ moves $-35~\mathrm{cm^{-1}}$ on
going from $\mathrm{TDX}$ to $\mathrm{TD}$. The couplet of the
\emph{two-chromophore} construct is therefore centred $35~\mathrm{cm^{-1}}$
--- $27\%$ of $\delta$ --- away from that construct's own absorption maximum,
and in a fit whose only splitting freedom is symmetric about that centre, a
mis-centring of this size is absorbed by $\delta$ and by the amplitude
asymmetry.

\emph{(iii) The fit's own widths violate the common-lineshape assumption.} The
two couplet components are given independent widths, and the fitted values
differ by a factor $1.50$ ($293.5$ against $439.3~\mathrm{cm^{-1}}$). This is
assumption~4 below, broken inside the analysis that the coupling is extracted
from. It is also why the couplet looks strongly non-conservative in amplitude
($A_{1}/|A_{2}|=1.58$) while its \emph{areas} cancel to $5.4\%$ of the positive
lobe ($A\pi\gamma$: $3310$ against $3132$), area being the quantity
conservativity actually constrains.

\emph{(iv) No uncertainty is quoted on $\delta$.} The supplement describes a
$5000$-sample resampling procedure, and reports bootstrap-based comparisons for
the absorption decomposition, but Table~S3 gives $\delta$ to five significant
figures with no interval. In the overlapped limit $\delta\ll\gamma$ the couplet
shape depends on $\delta$ and the amplitude only through the product $A\delta$
(Eq.~\ref{eq:derivative}), so $\delta$ is fixed solely by the departure from
the derivative limit. At $2\delta/\gamma\approx0.6$--$0.9$ that departure is
present but small, and the fit is correspondingly ill-conditioned: the reported
precision is a property of the constraint, not of the data.

None of this makes the measurement wrong. It makes $\delta$ a
model-conditioned quantity --- the splitting a degenerate dimer with these
widths would need in order to reproduce the observed couplet --- which is
precisely $J_{\mathrm{app}}$ of Eq.~(\ref{eq:japp}), not $J$.

\subsection{A moment observable immune to detuning}
\label{sec:moment}

There is a reduction of the couplet that survives arbitrary overlap. Because
$g$ is normalised and common, and each couplet is conservative
($\int\Delta A\,d\nu=0$, which also makes the moment origin-independent), the
signed first moment of the \emph{entire} CD band is
\begin{equation}
M_{1}\equiv\int \nu\,\Delta A(\nu)\,d\nu
\;\propto\;\frac{1}{N}\sum_{i}R^{i}_{+}\left(\nu^{i}_{+}-\nu^{i}_{-}\right)
=\frac{1}{N}\sum_{i}R^{i}_{+}\,\Omega_{i}
=\pm\,\frac{\pi\bar\nu}{N}\sum_{i}|J_{i}T_{i}| ,
\label{eq:moment}
\end{equation}
where the last step uses the sum rule. \emph{All dependence on the detuning
has cancelled}: $M_{1}$ measures the ensemble mean of $|JT|$ exactly, at any
degree of band overlap and for any distribution of $\bar\nu_i$. The practical
price is that extracting $|J|$ from $M_{1}$ requires the geometric factor $T$
(known here from structure), an absolute intensity calibration, and an
integration window free of vibronic and neighbouring-transition
contamination. Where those conditions can be met, Eq.~(\ref{eq:moment}) --- not
the peak separation --- is the model-robust way to reduce a couplet to a
coupling; related sum-rule analyses of aggregate CD are given in
Refs.~\cite{somsen1996,vanamerongen2000}.

The published dVenus fit cannot be reduced this way as it stands, and the
reason is instructive. $M_{1}$ is origin-independent only when the integrated
band vanishes, and for the Table~S3 parameters it does not: the couplet is
conservative in area to $5.4\%$, and the third component --- which the paper
fits at $20{,}774~\mathrm{cm^{-1}}$ but never assigns --- carries $18\%$ of the
couplet's area with the same sign as the negative lobe. Until that component is
assigned (its $\approx 1450~\mathrm{cm^{-1}}$ offset from the couplet centre is
a C=C/C=N stretch, so vibronic is the natural reading, but that is an inference
and not the authors') the integration window is not defined, and a first moment
taken over it inherits an origin-dependent contribution comparable to the
signal. A moment reduction therefore needs the \emph{raw} difference spectrum
and a defensible window, not the published decomposition. (For Lorentzian
components $\int\nu\,L\,d\nu$ exists only as a principal value; that is
harmless once the zeroth moment vanishes, and fatal when it does not.)

\section{What a degenerate-model fit returns}

An analyst who assumes the two chromophores are isoenergetic sets $\Delta=0$ in
Eq.~(\ref{eq:omega}), so that the measured component separation is interpreted
as $2|J|$. In the resolved regime the recovered value is therefore
\begin{equation}
\boxed{\;
J_{\mathrm{app}}=\tfrac{1}{2}\langle\Omega\rangle_{\mathrm{CD}}
=\frac{1}{2}\,\frac{\sum_{i} |J_{i}T_{i}|}{\sum_{i} |J_{i}T_{i}|/\Omega_{i}} \;}
\label{eq:japp}
\end{equation}
For a single configuration this is simply
$J_{\mathrm{app}}=\tfrac{1}{2}\sqrt{\Delta^{2}+4J^{2}}\ge |J|$, with equality
only at $\Delta=0$: within the static electronic model the inferred coupling
can only ever be too large, never too small. Taking $J_{i}=J>0$ and $T_{i}=T$
fixed, the inflation factor is
\begin{equation}
\frac{J_{\mathrm{app}}}{J}
=\frac{\langle\Omega\rangle_{\mathrm{CD}}}{2J}
\;\xrightarrow[\;|\Delta|\gg 2|J|\;]{}\;
\frac{\langle|\Delta|\rangle_{\mathrm{CD}}}{2J},
\end{equation}
i.e.\ in the strongly detuned limit the apparent coupling measures the
\emph{detuning}, not the coupling at all. In the overlapped regime of
Sec.~\ref{sec:overlap} the situation is worse still: peak reading returns
$J_{\mathrm{app}}\approx\sigma_g$, which is bounded below by the linewidth
however small $J$ and $\Delta$ may be.

\section{Assumptions, and which of them can fail}

\begin{enumerate}
\item \textbf{Coupled-oscillator CD.} Intrinsic monomer rotational strength,
inter-chromophore electric--magnetic cross terms, and quadrupole contributions
are neglected \cite{tinoco1962,jurinovich2014}. This is the standard
exciton-chirality treatment and is appropriate here because the experimental
observable is a \emph{difference} spectrum, constructed precisely to cancel
the monomer contribution --- with the caveat that the cancellation is exact
only if dimerisation leaves the intrinsic monomer CD unchanged.

\item \textbf{Sign coherence.} $\mathrm{sign}(J_{i}T_{i})$ is taken to be
common to all configurations. If the coupling (or the chirality sense) changes
sign across the ensemble, sub-populations contribute couplets of opposite
sense that partially cancel; the weights in Eq.~(\ref{eq:weighted}) then
become signed and the moment analysis can fail outright (the denominator can
pass through zero). For a rigid dimer with $|J|$ well above its fluctuation
this is safe, but it must be checked on the trajectory.

\item \textbf{Static ensemble.} Each configuration is assumed frozen on the
optical timescale, so the spectrum is an incoherent sum,
Eq.~(\ref{eq:ensemble}). This requires the site-energy-gap correlation time to
exceed the dephasing time; in the opposite limit the ensemble is exchange
narrowed \cite{knapp1984} and the effective detuning spread collapses toward
its mean. Note the direction of this failure: narrowing drives
$J_{\mathrm{app}}$ \emph{down toward} $|J|$ (for $\langle\Delta\rangle=0$), so
violating this assumption reduces the inflation but cannot turn it into an
underestimate. \emph{This assumption must be checked, not assumed.} For these
constructs there is now a measured number to check it against: $128.0\pm8.1$~fs
for $\mathrm{dVenus\mbox{-}TD}$ by two-photon NSOM on a thin film (Table~S4 of
Ref.~\cite{nguyen2025}), against which the site-energy gap correlation time
from simulation must be compared --- with the caveat that a dried film is not
the solution sample the CD was measured on.

\item \textbf{Common lineshape.} Eq.~(\ref{eq:weighted}) uses the first
moment, which is exact provided $g$ is identical for every configuration and
for both bands. In reality the upper exciton band is lifetime-broadened by
downhill interexciton relaxation \cite{renger2002}, which correlates linewidth
with lobe and biases the moment separation. The published fit does not assume
it either: its two couplet components carry independent widths differing by a
factor $1.50$ (Sec.~\ref{sec:published}, point~iii).

\item \textbf{Two states per monomer.} Vibronic structure is excluded. This is
not only a fitting hazard --- and a live one here, since both published fits
require a third, unassigned band under the couplet, and the absorption
decomposition of the same spectra needs four Gaussians by AIC --- it is a
physical one. Vibronic coupling renormalises the
effective resonance coupling of the 0--0 region by a Franck--Condon factor,
$J_{\mathrm{eff}}\simeq\braket{0|0}^{2}J<J$ in the weak-exchange limit
\cite{simpson1957,hestand2018}, so a CD-derived splitting can also
\emph{under}-report the bare electronic coupling. This is the one mechanism in
this list that can reverse the sign of the bias.

\item \textbf{Resolved couplet, moment-returning fit.} Equations
(\ref{eq:weighted})--(\ref{eq:japp}) presuppose the resolved regime, and real
analyses fit component centres rather than moments. When the bands overlap the
fit is ill-conditioned ($J$, $\Delta$ and linewidth trade off against one
another), naive peak reading returns the linewidth
(Sec.~\ref{sec:overlap}), and only amplitude- or moment-based reductions
(Sec.~\ref{sec:moment}) retain coupling information.
\end{enumerate}

Assumptions 2, 3, 5 and 6 are the ones that can materially move the answer,
and none of them is guaranteed by the derivation itself.

\section{Consequence}

Three statements follow that do not depend on the numerical details:

\begin{itemize}
\item Within the static, purely electronic model, a coupling extracted from CD
under an assumption of degeneracy is an \emph{upper bound} on the true
coupling, configuration by configuration --- and if the couplet is unresolved,
peak reading inflates it further, to the linewidth scale.
\item The bound is model-bound, not unconditional. Motional narrowing only
tightens it toward $|J|$; but vibronic (Franck--Condon) renormalisation can
push the observed 0--0 splitting \emph{below} the bare electronic $2|J|$
\cite{simpson1957,hestand2018}. A CD-derived coupling is therefore neither a
reliable upper nor lower bound once vibronic structure is strong.
\item The discrepancy between a computed $J$ and a CD-derived ``coupling'' is
not by itself evidence that the computed $J$ is wrong. Establishing that
requires independent knowledge of $\Delta$ (from electronic structure --- a
single splitting measurement cannot separate $\Delta$ from $J$,
Eq.~(\ref{eq:omega})), or an observable that does not contain $\Delta$ at all,
such as the signed first moment of Eq.~(\ref{eq:moment}).
\end{itemize}

\subsection*{A second detuning-free observable, in absorption}

The CD moment of Sec.~\ref{sec:moment} is not the only reduction that cancels
$\Delta$. The exciton dipole strengths are
$|\muv_{\pm}|^{2}=\mu^{2}(1\pm\sin2\theta\cos\alpha)$ with $\alpha$ the angle
between the monomer dipoles, so the dipole-strength-weighted first moment of
the absorption band is
\begin{equation}
\bar{E}_{\mathrm{abs}}
=\bar{E}+\frac{\Omega}{2}\,\frac{2J}{\Omega}\cos\alpha
=\bar{E}+J\cos\alpha ,
\label{eq:absmoment}
\end{equation}
the $\Omega$ of the intensity asymmetry cancelling the $\Omega$ of the
splitting. Like Eq.~(\ref{eq:moment}) this survives averaging over any
distribution of $\Delta$; unlike it, it is carried by absorption, which is far
better determined than a difference-CD couplet. The same experiments supply it:
$-35~\mathrm{cm^{-1}}$ for the $\mathrm{TDX}\to\mathrm{TD}$ shift of the
$G_{1}$ absorption component \cite{nguyen2025}, and $-33.8~\mathrm{cm^{-1}}$
from the monomer-to-dimer shift $515.9\to516.8$~nm \cite{kim2019} --- two
independent measurements agreeing to $4\%$. It is a bound rather than a match,
because the environmental change on dimerisation (or on deleting a tyrosine)
shifts the site energies as well; but attributing the whole shift to the
exciton gives $|J|\le 35/|\cos\alpha|~\mathrm{cm^{-1}}$, which is
$35~\mathrm{cm^{-1}}$ for the head-to-tail geometry the authors invoke and
$53~\mathrm{cm^{-1}}$ at their own anisotropy-derived $|\cos\alpha|=0.660$.
Reconciling $J=130.79~\mathrm{cm^{-1}}$ with a $35~\mathrm{cm^{-1}}$ shift
requires $|\cos\alpha|\le0.27$, i.e.\ dipoles within $16^{\circ}$ of
perpendicular --- the opposite of the head-to-tail arrangement the same red
shift is cited as evidence for. The two geometric readings in the experimental
papers are already inconsistent with each other --- ``predominant head-to-tail''
from the red shift, $|\cos\alpha|=0.660$ from the limiting-anisotropy drop ---
so this tension is internal to the measurement, not a conflict with anything
computed here.

\section{Bypassing the inversion: CD directly from the QM/MM simulation}
\label{sec:direct}

Every difficulty above stems from \emph{inverting} the spectrum to a model
parameter. The inversion is unnecessary: the QM/MM machinery already in place
for site energies and couplings can produce the CD spectrum itself, which is
then compared with experiment with no degeneracy assumption, no lineshape
deconvolution, and no fit. Two complementary routes:

\subsection{Per-snapshot rotatory strengths (fit-free, exciton-model-free)}

For each MD snapshot, place \emph{both} chromophores (and any strongly
interacting residues) in the QM region, the remaining protein and solvent in
the (ideally polarizable) MM region, and compute the low-lying excited states
$k$ of the full dimer together with their rotatory strengths
\begin{equation}
R_{k}=\operatorname{Im}\big[\braket{0|\hat{\muv}|k}\cdot\braket{k|\hat{\mathbf{m}}|0}\big],
\end{equation}
evaluated in the velocity gauge (or with gauge-including orbitals) so the
result is origin-independent \cite{autschbach2009}. The ensemble CD spectrum
is then the direct sum
\begin{equation}
\Delta\varepsilon(\nu)\;\propto\;\nu\,
\Big\langle{\textstyle\sum_{k}}\,R_{k}\,g\big(\nu-\nu_{k}\big)\Big\rangle_{\mathrm{snapshots}},
\label{eq:directcd}
\end{equation}
with residual broadening $g$ representing only what the ensemble does not
already contain. Running the identical protocol on the monomer construct gives
the dimer-minus-monomer difference spectrum \emph{exactly as measured}. This
route automatically contains everything the exciton-chirality model discards:
intrinsic monomer CD, electric--magnetic cross terms, charge-transfer mixing,
and environment-induced chirality. It is the strategy of the
TDDFT/MMPol/PCM exciton-CD framework of Jurinovich \emph{et
al.}~\cite{jurinovich2014}; standard good-practice protocols for computed ECD
apply \cite{pescitelli2016}, and simplified TDDFT makes the per-snapshot cost
of hundreds of frames tractable \cite{bannwarth2014}. Any excited-state
method that yields transition electric \emph{and} magnetic dipoles (TDDFT,
sTD-DFT, EOM/STEOM response) can be substituted.

\subsection{Dynamic lineshape from the fluctuating exciton Hamiltonian}

The snapshot sum, like Eq.~(\ref{eq:ensemble}), is still a static
approximation (assumption~3). The QM/MM trajectory, however, already provides
the \emph{time series} $E_{A}(t)$, $E_{B}(t)$, $J(t)$ and the geometric
factors $\Rv_{AB}(t)$, $\muv_{a}(t)$; the CD lineshape follows from linear
response as the Fourier transform of the electric--magnetic cross-correlation
function. At the exciton level,
\begin{equation}
\Delta\varepsilon(\omega)\;\propto\;\omega\,
\operatorname{Re}\!\int_{0}^{\infty}\!dt\;e^{i\omega t}\,
\Big\langle{\textstyle\sum_{a\neq b}}\;
\pi\bar\nu\,\Rv_{ab}\cdot\big(\muv_{a}\times\muv_{b}\big)\,
\big[U(t)\big]_{ba}\Big\rangle_{\mathrm{traj}},
\label{eq:response}
\end{equation}
with the time-ordered propagator of the fluctuating Hamiltonian
\begin{equation}
U(t)=\mathcal{T}\exp\!\Big[-\tfrac{i}{\hbar}\!\int_{0}^{t}\!H(\tau)\,d\tau\Big],
\end{equation}
averaged over trajectory time origins, optionally with a second-order-cumulant
treatment of the fast intramolecular modes through their spectral density
\cite{renger2002,zuehlsdorff2019}. Equation~(\ref{eq:response}) interpolates
between the limits by itself: for slow fluctuations it reduces to the static
sum (\ref{eq:ensemble}); for fast fluctuations it produces the exchange-narrowed
spectrum \cite{knapp1984}. It therefore \emph{decides} assumption~3 rather
than assuming it, and it reuses precisely the site-energy and coupling
trajectories this project already generates. Disorder-averaged exciton CD of
pigment--protein aggregates computed in this spirit is standard in the
photosynthesis literature \cite{somsen1996,vanamerongen2000,renger2002}.

\begin{thebibliography}{99}

\bibitem{davydov1962}
A.~S. Davydov, \emph{A Theory of Molecular Excitons} (McGraw-Hill, New York,
1962).

\bibitem{kasha1965}
M.~Kasha, H.~R. Rawls, M.~A. El-Bayoumi,
The exciton model in molecular spectroscopy,
\emph{Pure Appl. Chem.} \textbf{11}, 371--392 (1965).

\bibitem{tinoco1962}
I.~Tinoco, Jr.,
Theoretical aspects of optical activity. Part two: polymers,
\emph{Adv. Chem. Phys.} \textbf{4}, 113--160 (1962).

\bibitem{schellman1975}
J.~A. Schellman,
Circular dichroism and optical rotation,
\emph{Chem. Rev.} \textbf{75}, 323--331 (1975).

\bibitem{harada1983}
N.~Harada, K.~Nakanishi, \emph{Circular Dichroic Spectroscopy: Exciton
Coupling in Organic Stereochemistry} (University Science Books, Mill Valley,
1983).

\bibitem{berova2007}
N.~Berova, L.~Di Bari, G.~Pescitelli,
Application of electronic circular dichroism in configurational and
conformational analysis of organic compounds,
\emph{Chem. Soc. Rev.} \textbf{36}, 914--931 (2007).

\bibitem{rodger1997}
A.~Rodger, B.~Nord\'en, \emph{Circular Dichroism and Linear Dichroism}
(Oxford University Press, Oxford, 1997).

\bibitem{somsen1996}
O.~J.~G. Somsen, R.~van Grondelle, H.~van Amerongen,
Spectral broadening of interacting pigments: polarized absorption by
photosynthetic proteins,
\emph{Biophys. J.} \textbf{71}, 1934--1951 (1996).

\bibitem{vanamerongen2000}
H.~van Amerongen, L.~Valkunas, R.~van Grondelle,
\emph{Photosynthetic Excitons} (World Scientific, Singapore, 2000).

\bibitem{knapp1984}
E.~W. Knapp,
Lineshapes of molecular aggregates: exchange narrowing and intersite
correlation,
\emph{Chem. Phys.} \textbf{85}, 73--82 (1984).

\bibitem{simpson1957}
W.~T. Simpson, D.~L. Peterson,
Coupling strength for resonance force transfer of electronic energy in
van der Waals solids,
\emph{J. Chem. Phys.} \textbf{26}, 588--593 (1957).

\bibitem{hestand2018}
N.~J. Hestand, F.~C. Spano,
Expanded theory of H- and J-molecular aggregates: the effects of vibronic
coupling and intermolecular charge transfer,
\emph{Chem. Rev.} \textbf{118}, 7069--7163 (2018).

\bibitem{renger2002}
T.~Renger, R.~A. Marcus,
On the relation of protein dynamics and exciton relaxation in
pigment--protein complexes: an estimation of the spectral density and a
theory for the calculation of optical spectra,
\emph{J. Chem. Phys.} \textbf{116}, 9997--10019 (2002).

\bibitem{jurinovich2014}
S.~Jurinovich, G.~Pescitelli, L.~Di Bari, B.~Mennucci,
A TDDFT/MMPol/PCM model for the simulation of exciton-coupled circular
dichroism spectra,
\emph{Phys. Chem. Chem. Phys.} \textbf{16}, 16407--16418 (2014).

\bibitem{autschbach2009}
J.~Autschbach,
Computing chiroptical properties with first-principles theoretical
chemistry: background and illustrative examples,
\emph{Chirality} \textbf{21}, E116--E152 (2009).

\bibitem{pescitelli2016}
G.~Pescitelli, T.~Bruhn,
Good computational practice in the assignment of absolute configurations by
TDDFT calculations of ECD spectra,
\emph{Chirality} \textbf{28}, 466--474 (2016).

\bibitem{bannwarth2014}
C.~Bannwarth, S.~Grimme,
A simplified time-dependent density functional theory approach for
electronic ultraviolet and circular dichroism spectra of very large
molecules,
\emph{Comput. Theor. Chem.} \textbf{1040--1041}, 45--53 (2014).

\bibitem{zuehlsdorff2019}
T.~J. Zuehlsdorff, A.~Montoya-Castillo, J.~A. Napoli, T.~E. Markland,
C.~M. Isborn,
Optical spectra in the condensed phase: capturing anharmonic and vibronic
features using dynamic and static approaches,
\emph{J. Chem. Phys.} \textbf{151}, 074111 (2019).

\bibitem{kim2019}
Y.~Kim, H.~L. Puhl, E.~Chen, G.~H. Taumoefolau, T.~A. Nguyen, D.~S. Kliger,
P.~S. Blank, S.~S. Vogel,
Venus$_{\mathrm{A206}}$ dimers behave coherently at room temperature,
\emph{Biophys. J.} \textbf{116}, 1918--1930 (2019).

\bibitem{nguyen2025}
T.~Nguyen, H.~L. Puhl, E.~Chen, K.~Hines, C.~Rani, P.~S. Blank, B.~Carlotti,
Y.~Kim, T.~Goodson, S.~S. Vogel,
Anomalous photophysical behaviors attributed to excitonic coupling in
fluorescent proteins,
\emph{Biophys. J.} \textbf{124}, 4293--4309 (2025).

\end{thebibliography}

\end{document}
